\scanpage{365}\mspage{281}
4.13. \uline{\deleted{\unsure{Parties}}\unsure{Ensembles} accessibles d'objets d'une
catégorie accessible}.

\uline{Définition} 4.13.1. Soient $M$ une catégorie
accessible, $\deleted{M}\unsure{S|\mathscr S}$ une partie de $\operatorname{Ob}M$,
$\underline{\deleted{M}S}$ la sous-catégorie pleine de $M$ définie
par $M$. On dit que $\deleted{M}S$ est une
\uline{partie accessible} de $\operatorname{Ob}M$, ou encore
que $\underline{\deleted{M}S}$ est une \uline{\inserted{ss-}catégorie} ``\uline{accessible
dans $M$}'', si a) $\underline{\deleted{M}S}$ est \unsure{strictemt|strictmt|strictement} pleine,
\unsure{(i.e.|ie} pour \deleted{$x\in M$, et $y\in\operatorname{Ob}M$,} $x,y$
deux objets isomorphes de $M$, $x\unsure{+|\in}\deleted{M}S$
implique $y\in\deleted{M}S$), \deleted{et} b) $\underline{\deleted{M}S}$ est
accessible, c) le foncteur d'inclusion
$\underline{\deleted{M}S}\unsure{-|\longrightarrow}M$ est accessible.

Soit\deleted{\unsure{ent}} $\pi_0$ un cardinal utile pour
$M$, i.e. tel que $M\unsure{\underline{\approx}|\simeq}\mathrm{Ind}_{\pi_0}(M_{\pi_0})$,
soit $\pi_1=\operatorname{Sup}(\pi_0,\operatorname{Card\,red} M_{\pi_0})$. Soit
d'autre part $\pi'_0$ un cardinal tel que
% archive sheet number: 364
